\chapter{Amylase}

\section*{What Is This Test?}
Amylase is an enzyme that catalyzes the hydrolysis of complex carbohydrates (starch and glycogen) into smaller sugars (maltose, maltotriose, and limit dextrins) by cleaving alpha-1,4-glycosidic bonds. In humans, amylase is produced primarily by two organs: the pancreas (pancreatic amylase, P-type isoenzyme, approximately 40--50\% of serum amylase) and the salivary glands (salivary amylase, S-type isoenzyme, approximately 50--60\% of serum amylase). Small amounts are also produced by the ovaries, fallopian tubes, testes, and small intestine. Pancreatic amylase is secreted into the duodenum via the pancreatic duct, where it initiates starch digestion. Salivary amylase is secreted into the mouth, beginning carbohydrate digestion. Amylase is a relatively small enzyme (molecular weight ~55,000 Da), and approximately 25\% is filtered by the glomerulus and excreted in urine (amylasuria). In acute pancreatitis, amylase leaks from damaged acinar cells into the blood. Serum amylase rises within 2--12 hours of symptom onset, peaks at 12--72 hours, and returns to normal within 3--5 days (serum half-life approximately 2 hours). Traditionally, amylase was the primary diagnostic test for acute pancreatitis, but it has been largely superseded by lipase, which offers superior sensitivity (85--100\% vs. 70--95\%) and specificity, a longer diagnostic window (lipase stays elevated for 8--14 days), and is not affected by non-pancreatic conditions that elevate amylase (salivary disease, macroamylasemia). Current guidelines from the American College of Gastroenterology (ACG) and the International Association of Pancreatology recommend lipase over amylase for the diagnosis of acute pancreatitis \cite{tenner2013american,banks2013classification}.

\section*{Alternative Names}
\begin{itemize}
  \item Serum Amylase
  \item Blood Amylase
  \item Amylase, Total
  \item Pancreatic Amylase
  \item Salivary Amylase
  \item Urine Amylase
  \item Amylase/Creatinine Clearance Ratio
\end{itemize}
\section*{Principle}
Amylase is measured by enzymatic methods using synthetic substrates. The most common method uses 2-chloro-4-nitrophenyl-alpha-D-maltotrioside (CNP-G3) or ethylidene-blocked p-nitrophenyl-alpha-D-maltoheptaoside (EPS-G7). The substrate is cleaved by alpha-amylase, releasing smaller fragments. The auxiliary enzyme alpha-glucosidase hydrolyzes these fragments, ultimately releasing the chromophore 2-chloro-4-nitrophenol (CNP) or p-nitrophenol, which is yellow and measured spectrophotometrically at 405 nm. The rate of color production is proportional to amylase activity. Some methods use monoclonal antibodies to specifically inhibit salivary amylase, allowing measurement of pancreatic amylase alone. Urine amylase (often measured as the amylase/creatinine clearance ratio) is historically used to differentiate macroamylasemia (normal clearance ratio <1\%) from pancreatitis (elevated clearance ratio >3\%), though this test is now rarely performed. This test is NOT performed by hematology analyzers.

\section*{History}
Amylase was the first enzyme to be discovered, identified in 1833 by Anselme Payen and Jean-Fran\c{c}ois Persoz \cite{payen1833diastase}. The association between serum amylase and acute pancreatitis was first reported by Elman in 1929 \cite{elman1929blood}. Amylase became the standard diagnostic test for pancreatitis throughout the 20th century. The isoenzyme fractionation into pancreatic (P-type) and salivary (S-type) forms was developed in the 1970s. Lipase began to replace amylase as the preferred test in the 1990s-2000s, and current ACG guidelines (2013) recommend lipase as the diagnostic test of choice for acute pancreatitis \cite{tenner2013american}.

\section*{Why This Test Is Ordered}
\begin{itemize}
  \item Adjunctive diagnosis of acute pancreatitis: amylase >3$\times$ ULN in the context of acute epigastric abdominal pain radiating to the back, with nausea/vomiting. Lipase is preferred and should always be ordered simultaneously
  \item Evaluation of suspected chronic pancreatitis: amylase is often normal in chronic pancreatitis (acinar cell destruction reduces enzyme production) and is not useful for diagnosis
  \item Detection of pancreatic duct obstruction or pancreatic pseudocyst: amylase may be persistently mildly elevated
  \item Evaluation of macroamylasemia: persistently elevated total amylase without clinical pancreatitis, caused by amylase complexed with immunoglobulins (usually IgA or IgG), which prevents renal clearance. This is a benign condition
  \item Diagnosis of salivary gland disorders: parotitis (mumps), sialadenitis, Sj\"{o}gren's syndrome, salivary duct obstruction (stones)
  \item Monitoring of post-ERCP pancreatitis: amylase checked 2--4 hours after ERCP to detect early pancreatic injury
\end{itemize}

\section*{Primary vs Secondary Test}
Amylase is now a secondary test for pancreatitis, with lipase being the preferred primary marker per ACG guidelines. Amylase retains utility for specific non-pancreatic indications and when lipase is unavailable.

\section*{How This Test Is Performed}
A healthcare professional draws blood from a vein into a serum separator tube (gold-top) or lithium heparin tube (green-top). EDTA or citrate tubes should NOT be used because they chelate calcium, which is a required cofactor for amylase activity. Hemolysis should be avoided. The sample is centrifuged and analyzed on an automated chemistry analyzer. Results are available within 1--2 hours. The tube must be capped to avoid contamination.

\section*{What Preparation Is Needed}
No fasting is required, but ethanol intake should be avoided for 24 hours before testing (alcohol can mildly elevate amylase). Inform the provider of: (1) Medications that can elevate amylase: opiates (morphine, codeine---cause spasm of the sphincter of Oddi, elevating amylase even in the absence of pancreatitis), aspirin, corticosteroids, furosemide, thiazide diuretics, valproic acid, azathioprine, and didanosine. (2) Recent ERCP (amylase routinely checked post-procedure). (3) Chronic pancreatitis: amylase may be normal or low despite disease due to pancreatic acinar destruction. (4) A normal amylase does NOT exclude acute pancreatitis---up to 20--30\% of patients with acute pancreatitis have normal amylase, especially in alcoholic pancreatitis (repeat insults deplete acinar enzyme stores), hypertriglyceridemia-induced pancreatitis (triglycerides interfere with the assay), and late presentation (>48--72 hours after symptom onset when amylase has already normalized).

\section*{Contraindications and Precautions}
\begin{itemize}
  \item No absolute contraindications. Critical considerations:
  \item (1) Amylase is less sensitive (70--95\%) and less specific than lipase for acute pancreatitis. If there is clinical suspicion for pancreatitis, lipase should always be ordered.
  \item (2) Amylase has a short half-life (approximately 2 hours) and normalizes within 3--5 days. Patients presenting >48--72 hours after symptom onset may have normal amylase despite ongoing pancreatitis. Lipase remains elevated for 8--14 days.
  \item (3) Non-pancreatic conditions causing elevated amylase: salivary gland disease (parotitis from mumps, sialolithiasis, Sj\"{o}gren's, radiation injury), intestinal obstruction, perforated peptic ulcer, acute cholecystitis, ruptured ectopic pregnancy, ovarian cyst/tumor, renal failure (decreased amylase clearance), macroamylasemia, diabetic ketoacidosis, anorexia nervosa/bulimia (salivary amylase), and medications (opiates cause sphincter of Oddi spasm).
  \item (4) Hypertriglyceridemia: triglycerides >500 mg/dL can cause spuriously normal amylase because the turbidity and lipid interfere with the spectrophotometric assay. Lipase is also affected but to a lesser degree. Dilution or ultracentrifugation can resolve interference.
  \item (5) Macroamylasemia: amylase complexed with immunoglobulin (IgA or IgG) cannot be filtered by the kidney, causing persistently elevated serum amylase with low urine amylase. The amylase-to-creatinine clearance ratio (ACCR) <1\% (normal 2--5\%) is diagnostic. Benign condition requiring no treatment.
  \item (6) ED diagnostic protocols: lipase is now the standard, and amylase may not need to be ordered in many institutions.
\end{itemize}
\section*{Risks}
Minimal risk from venipuncture.

\section*{What Happens After the Test}
Results available within 1--2 hours. For suspected acute pancreatitis, amylase >3$\times$ ULN (often >200--300 IU/L) with characteristic epigastric pain supports the diagnosis. Two of the following three criteria are required for diagnosis: (1) characteristic abdominal pain (acute onset, severe, epigastric, radiating to back), (2) serum amylase and/or lipase >3$\times$ ULN, (3) characteristic findings on abdominal imaging (CT, MRI, or ultrasound). If amylase is elevated but lipase is normal, consider non-pancreatic causes (salivary, macroamylasemia, renal failure). If both amylase and lipase are normal, but clinical suspicion for pancreatitis is high, obtain abdominal CT with IV contrast to evaluate for pancreatitis, complications (necrosis, pseudocyst, abscess), and alternative diagnoses. In acute pancreatitis, Ranson criteria (at 48 hours) and BISAP score (at 24 hours) are used to predict severity and mortality.

\section*{Understanding Results}

\subsection*{Below Normal}
\begin{itemize}
  \item Low amylase is not a common clinical concern but may indicate chronic pancreatitis with extensive acinar cell destruction (``burned-out'' pancreas), resulting in low or normal amylase despite disease.
  \item Cystic fibrosis: pancreatic insufficiency leads to low amylase production.
  \item Advanced chronic liver disease.
  \item Post-pancreatectomy.
  \item Normal amylase does not exclude acute pancreatitis (up to 20--30\% of cases).
\end{itemize}

\subsection*{Above Normal}
\begin{itemize}
  \item \textbf{Acute pancreatitis (most important):} Amylase >3$\times$ ULN (typically 200--1,000+ IU/L). Common causes: gallstones (40\%), alcohol (30\%), idiopathic (15\%), hypertriglyceridemia (>1,000 mg/dL), post-ERCP, medications (azathioprine, valproic acid, didanosine, sulfonamides, thiazides), hypercalcemia, trauma, pancreatic cancer, autoimmune pancreatitis, scorpion stings, and genetic (hereditary pancreatitis---PRSS1 mutation).
  \item \textbf{Salivary gland disease:} Mumps (paramyxovirus parotitis---markedly elevated salivary amylase), sialadenitis (bacterial or viral), sialolithiasis (salivary duct stone), Sj\"{o}gren's syndrome, radiation injury to salivary glands, anorexia nervosa/bulimia (vomiting-induced sialadenosis).
  \item \textbf{Gastrointestinal:} Intestinal obstruction, perforated peptic ulcer (amylase absorbed from gut lumen across damaged mucosa), acute cholecystitis (mild elevation), mesenteric ischemia, acute appendicitis, ruptured ectopic pregnancy, ovarian cysts and tumors.
  \item \textbf{Macroamylasemia:} Amylase complexed with immunoglobulins (IgA or IgG); persistent benign hyperamylasemia; low urine amylase and ACCR <1\%. PEG precipitation confirms.
  \item \textbf{Renal failure:} Decreased glomerular filtration reduces amylase clearance, causing mild-to-moderate elevation. Both amylase and lipase may be elevated in CKD (amylase 2--3$\times$ ULN, lipase 1.5--2$\times$ ULN in ESRD).
  \item \textbf{Medications:} Opiates (morphine, codeine---sphincter of Oddi spasm), aspirin, corticosteroids, thiazides, furosemide, valproic acid, azathioprine, didanosine, alcohol.
\end{itemize}

\section*{Normal Results}
\begin{itemize}
  \item \textbf{Total amylase:} 30--110 IU/L (varies by method; typical 40--140 IU/L)
  \item \textbf{Pancreatic amylase:} 13--53 IU/L (typically 40--50\% of total)
  \item \textbf{Salivary amylase:} 17--57 IU/L (typically 50--60\% of total)
  \item \textbf{Amylase >3$\times$ ULN:} Diagnostic threshold for acute pancreatitis (with clinical context)
  \item \textbf{Urine amylase:} 2--35 IU/h (24-hour collection). Amylase-to-creatinine clearance ratio (ACCR): 2--5\%
  \item \textbf{Amylase peak in pancreatitis:} Rises 2--12 hours; peaks 12--72 hours; normalizes 3--5 days
  \item \textbf{Amylase half-life:} Approximately 2 hours
\end{itemize}

\section*{Follow-up Tests}
\begin{itemize}
  \item \textbf{Lipase:} Always ordered with amylase for pancreatitis evaluation. More sensitive (85--100\%) and specific; preferred test per ACG guidelines.
  \item \textbf{Abdominal CT with IV contrast:} Gold standard imaging for pancreatitis diagnosis, severity assessment (Balthazar score), and detection of complications (necrosis, pseudocyst, abscess, vascular complications).
  \item \textbf{Right upper quadrant ultrasound:} First-line for gallstone pancreatitis (most common cause).
  \item \textbf{Triglyceride level:} If hypertriglyceridemia-induced pancreatitis suspected (>1,000 mg/dL).
  \item \textbf{Amylase isoenzymes:} If macroamylasemia or salivary source is suspected.
  \item \textbf{Amylase/creatinine clearance ratio (ACCR):} <1\% confirms macroamylasemia.
  \item \textbf{Pancreatic protocol CT or MRCP:} For suspected pancreatic mass, ductal abnormalities, or chronic pancreatitis.
\end{itemize}

\section*{References}
\bibliographystyle{unsrtnat}
\bibliography{references}

\clearpage
\section*{Regulatory Status and Authorization Date}
This chapter describes a test category, not a specific commercial product. FDA, CDSCO, EU IVDR, and MHRA approval or registration dates therefore cannot be assigned without the assay manufacturer, product name, instrument, and intended use. For laboratory-developed tests, an individual agency approval date may not exist. See the Preface for the interpretation of regulatory status.

\clearpage
